\section{Abstract}
Streaming Speech Translation (SST) requires a delicate balance between real-time responsiveness and translation accuracy. Current simultaneous translation methods often suffer from ``early commitment'' errors, where a model generates a final translation before sufficient audio context is available. We propose a framework for \textbf{Selective Test-Time Reasoning (STTR)}, which dynamically allocates additional inference compute only when the model identifies high uncertainty in its own predictions. By focusing refinement steps on linguistically complex or ambiguous audio prefixes, we aim to improve the quality-latency Pareto frontier. We will evaluate our approach on the MuST-C and CoVoST2 benchmarks, measuring performance through BLEU and Length-Adaptive Average Lagging (LAAL).

% \section{Introduction and Motivation}
% Streaming translation is fundamentally constrained by the sequential arrival of input. Unlike offline systems, SST models must decide when to commit to an output word while possessing only partial information. This leads to critical failures in handling structural ambiguities, such as verb-final constructions in German, or high-stakes content like negations and numerical units. Once a token is committed to the user in a real-time stream, it cannot be easily revised, making the initial decision-making process high-risk.

% Recent advancements in Large Language Models (LLMs) suggest that additional test-time computation through self-correction or multi-pass generation can resolve complex reasoning tasks. However, applying these techniques uniformly in a streaming context would lead to unacceptable latency spikes. Our work is motivated by the insight that not all audio segments are equally difficult. By implementing a selective gating mechanism, we can maintain the low latency of standard models for simple prefixes while utilizing deeper reasoning only for the subset of inputs that require it.


\section{Introduction and Motivation}
Streaming translation is fundamentally constrained by the sequential arrival of input. Unlike offline systems that process complete utterances, SST models must decide when to commit to an output word while possessing only partial information \citep{ma2019stacl}. This leads to what is known as the ``early commitment'' problem, resulting in critical failures when handling structural ambiguities, such as verb-final constructions in German or SOV languages \citep{gu2017learning}. Furthermore, errors in high-stakes content like negation markers, numerical units, and named entities alter meaning and are difficult for streaming models to correct once the prefix is emitted \citep{papi2025real}.

Recent advancements in Large Language Models (LLMs) demonstrate that additional test-time computation through self-correction \citep{madaan2023selfrefine} or chain-of-thought reasoning \citep{wei2022cot} can significantly resolve complex reasoning tasks. \citet{snell2024scaling} further show that optimally scaling inference-time compute can sometimes be more effective than increasing model parameters. However, applying these techniques uniformly in a streaming context would lead to unacceptable latency spikes. Our work is motivated by the insight that not all audio segments are equally difficult. By implementing a selective gating mechanism, we can maintain the low latency of standard models for simple prefixes while utilizing deeper reasoning only for the subset of inputs that exhibit high semantic or structural uncertainty \citep{schuster2022calm}.

\section{Related Work}
Current research in SST mainly focuses on fixed policies like \textit{wait-k} \citep{ma2019stacl} or monotonic attention \citep{ma2020mma}. While these establish a reliable latency backbone, they lack the flexibility to adapt to varying input difficulty. In the offline domain, refinement techniques such as deliberation networks \citep{xia2017deliberation} and mask-predict \citep{ghazvininejad2019maskpredict} have shown quality improvements, but their high computational cost has limited their adoption in real-time pipelines \citep{dou2024iterative}.

The emerging field of test-time compute scaling \citep{snell2024scaling, wang2023selfconsistency} provides a new paradigm for enhancing model performance without increasing training costs. Notably, \citet{li2025tts_mt} demonstrated that domain-specific fine-tuning is essential for machine translation models to benefit from reasoning steps. Our work fills a critical gap by integrating adaptive compute logic, inspired by CALM \citep{schuster2022calm}, with uncertainty estimation \citep{fomicheva2020unsupervised, cheng2024entropy} to create the first difficulty-aware reasoning framework for streaming speech-to-text translation.

\section{Proposed Methodology}
We treat the translation process as a sequential decision task where each commit point is evaluated for potential refinement. For our baseline, we utilize a standard \textit{wait-k} policy with beam search. Building upon this, we introduce our primary contribution: Uncertainty-Triggered Refinement (UTR).

In the UTR pipeline, the model first generates a draft translation prefix. We then calculate a confidence score based on similarity-sensitive entropy. If the uncertainty is below a calibrated threshold $\tau$, the system commits the draft immediately to minimize latency. If the uncertainty exceeds $\tau$, the system triggers a lightweight second-pass refinement, either through multi-candidate re-ranking or iterative masking to correct potential early commitment errors. This selective application ensures that the system only incurs additional latency when the probability of a quality gain is high.

\section{Experimental Design}
Our experiments will be conducted on the MuST-C (English-German) and CoVoST2 datasets. English-German is specifically chosen as the primary test bed due to its significant word-order divergence, which frequently triggers early commitment errors in streaming models.

We will evaluate the system using a dual-metric approach. Translation quality will be measured via BLEU and COMET, while latency will be quantified using Average Lagging (AL) and LAAL via the SimulEval toolkit. To demonstrate the efficiency of our selective approach, we will compare it against two critical baselines: an ``Always-Refine'' system that serves as a quality upper bound, and a ``Larger-Beam'' search baseline that uses an equivalent total compute budget. This will allow us to isolate the benefits of adaptive reasoning from simple increased search breadth.

\section{Depth and Ablation Studies}
To understand the mechanics of selective reasoning, we will conduct a series of ablation studies. We will sweep the uncertainty threshold $\tau$ to trace the quality-latency Pareto frontier, identifying the optimal balance for various real-time use cases. Furthermore, we will compare different uncertainty signals, such as token-level entropy versus predictive variance, to determine which best correlates with translation failures. Finally, a linguistically-driven error analysis will be performed to verify if our method reduces errors in negation, numerals, and long-distance reordering, which are the most frequent failure points in streaming systems.

\section{Risks and Timeline}
The primary risk involves the potential mismatch between model uncertainty and actual translation errors; a model might be ``confidently wrong,'' causing the gate to fail \citep{wang2020calibration}. We mitigate this by calibrating our uncertainty metrics on a held-out development set. 

\textbf{Timeline:} Week 1: Reproduce wait-$k$ baseline and SimulEval pipeline. Week 2: Implement uncertainty estimation. Week 3: Integrate selective refinement and tradeoff experiments. Week 4: Ablation studies and report writing.

\section{Team and Resources}
The project will be hosted on GitHub, utilizing the Fairseq and SimulEval frameworks. Key references include the STACL framework \citep{ma2019stacl} for simultaneous policies and the MuST-C corpus \citep{digangi2019mustc} for benchmarking. For the division of labor, Jeng-Yue will take initial ownership of the data pipeline and baseline modeling. Wilson will support the core model implementation, while Haoling will lead the ablation studies and serve as the project manager.